% ============================================================================
\section{Parámetros numéricos}
\label{ap:parametros}
% ============================================================================

Todas las constantes del núcleo, con su valor y la razón que lo fija. Las
marcadas con $\ast$ se leen del entorno y se pueden cambiar sin modificar el
código.

\begin{table}[H]
\centering\small
\begin{tabular}{@{}L{3.5cm}L{3.9cm}L{7.6cm}@{}}
\toprule
Parámetro & Valor & Justificación \\
\midrule
\multicolumn{3}{@{}l}{\emph{Percepción}}\\
Tope de páginas & 40 & Justo por encima del PDF más largo del corpus (39 páginas). \\
Resolución de rasterizado & 150 DPI & Basta para leer una tabla como tabla; más resolución encarece la llamada sin ganancia. \\
Temperatura & 0 & Decodificación voraz: menor variabilidad entre lecturas. \\
Tiempo límite por llamada & 40 s & Una página típica tarda unos 5 s; 40 s ya es una llamada colgada. \\
Reintentos & 1 (solo de transporte) & Ante 429, 5xx o falta de respuesta; no ante 403 ni respuesta mal formada. \\
Espera entre intentos & 3 s & Peor caso por página: 83 s. \\
Modelo multimodal$^\ast$ & configurable & Puerto con una sola implementación; cambiarlo no toca las reglas. \\
\midrule
\multicolumn{3}{@{}l}{\emph{Resolución}}\\
Longitud mínima para corregir & 5 & Tokens más cortos tienen demasiados vecinos a distancia 1. \\
Distancia máxima & 1 (OSA) & Un error de dedo; la transposición cuenta como uno. \\
Prefijo de «misma cosa» & 5 caracteres & Sobre singulares; prefijos más cortos emparejan «lápida». \\
Palabras de la consulta & 6 & Los adjetivos discriminan; más palabras empeoran la búsqueda. \\
$\alpha$ (piso del ancla) & 0,6 & Acertar la cabeza basta para reconocer, no para autoaceptar. \\
$\beta$ (peso del Jaccard) & 0,4 & Con $\alpha$, el ancla perfecta vale 1. \\
$\gamma$ (aporte del coseno) & 0,15 & Pequeño: desempata, no decide. \\
$\theta$ (umbral de rescate) & 0,86 & Separa «la misma cosa dicha de otro modo» de «un tema parecido». \\
$\phi$ (piso del rescate) & 0,55 & Apenas sobre el umbral de reconocimiento. \\
Umbral de reconocimiento & 0,5 & Por debajo, \textsc{SinMatch}. \\
$\tau$ (autoaceptación)$^\ast$ & 0,75 & Por debajo, revisión humana; se puede subir cuando el catálogo esté más completo. \\
Opciones por renglón & $\leq3$ & Una por nivel; acota $\card{\Esp}\leq3^n$. \\
Dimensión de los vectores & 1024 & Fijada por la columna y el índice de la base. \\
Lote de \emph{embeddings} & 100 textos & Una petición por lote; evita cuerpos gigantes. \\
\bottomrule
\end{tabular}
\end{table}

\begin{table}[H]
\centering\small
\begin{tabular}{@{}L{3.5cm}L{3.9cm}L{7.6cm}@{}}
\toprule
Parámetro & Valor & Justificación \\
\midrule
\multicolumn{3}{@{}l}{\emph{Catálogo}}\\
Rango de precio en tiendas & \$1.000 -- \$400.000 & Descarta accesorios sueltos y lotes institucionales. \\
Propuestas por renglón & 3 (cuantiles 1/6, 1/2, 5/6) & Quien revisa compara pocas; se evitan los extremos. \\
Longitudes válidas de EAN & 8, 12, 13 & Las de 12 se normalizan a 13. \\
\midrule
\multicolumn{3}{@{}l}{\emph{Optimización}}\\
Límite de enumeración & $4096=2^{12}$ & Por debajo, el frente exacto cuesta milisegundos. \\
Población $N$ & $\min(200,\max(40,12n))$, par & Escala con las líneas decidibles. \\
Generaciones $G$ & $\min(120,\max(30,$ $20\log_{10}\card{\Esp}))$ & Escala con el logaritmo del espacio. \\
SBX & $p_c=0{,}9$, $\eta_c=15$ & Valores estándar \cite{deb1995sbx}. \\
PM & $p=1$ por individuo; $1/n$ por variable; $\eta_m=20$ & Valores estándar \cite{deb1996pm}. \\
Semilla & 0 & Reproducibilidad byte a byte. \\
Umbral de sobreprecio para Premium & $\sigma\geq0{,}10$ & Encuesta: 42 de 60 pagan como máximo un 5\,\% más. \\
Tope sin presupuesto (A~tu~Medida) & mitad del camino & Entre el total del Ahorrador y el del Premium. \\
\bottomrule
\end{tabular}
\end{table}
